\documentclass[UTF8,a4paper,10.5pt]{ctexart}
\usepackage[left=1.9cm,right=1.9cm,top=1.9cm,bottom=2.0cm]{geometry}
\usepackage{amsmath,amssymb}
\usepackage{enumitem}
\usepackage{fancyhdr}
\usepackage{lastpage}
\usepackage{listings}
\usepackage{xcolor}
\usepackage{array}

\pagestyle{fancy}
\fancyhf{}
\renewcommand{\headrulewidth}{0pt}
\fancyfoot[C]{第~\thepage~页\quad 共~\pageref{LastPage}~页}

\lstset{
    language=C,
    basicstyle=\ttfamily\small,
    keywordstyle=\color{blue!70!black}\bfseries,
    commentstyle=\color{green!45!black},
    stringstyle=\color{red!60!black},
    columns=fullflexible,
    keepspaces=true,
    breaklines=true,
    showstringspaces=false,
    aboveskip=3pt,
    belowskip=3pt
}

\newcommand{\ans}[1]{\hfill(\makebox[1.8cm][c]{#1})}
\newcommand{\fillans}[1]{\underline{\textbf{\ #1\ }}}

\begin{document}

\begin{center}
    {\Large\bfseries 2018 级信息学院《C 语言程序设计》考试试题（A）答案}
\end{center}

\noindent
\small
院系：\underline{\hspace{2.4cm}}\quad
班级：\underline{\hspace{2.4cm}}\quad
姓名：\underline{\hspace{2.2cm}}\quad
学号：\underline{\hspace{2.6cm}}\quad
成绩：\underline{\hspace{1.8cm}}
\normalsize

\vspace{6pt}
\begin{center}
\renewcommand{\arraystretch}{1.4}
\begin{tabular}{|c|*{7}{>{\centering\arraybackslash}p{1.4cm}|}}
\hline
题号 & 1 & 2 & 3 & 4 & 5 & 6 & 总分 \\
\hline
得分 &   &   &   &   &   &   &      \\
\hline
\end{tabular}
\end{center}

\section*{一、判断下列概念、语句或程序段的对错。（“$\times$”表示错，“$\checkmark$”表示对）（10 分）}

\begin{enumerate}[label=(\arabic*), leftmargin=2.2em, itemsep=4pt]
    \item \texttt{int x=0xabcd ; \quad // 对一个整型变量初始化} \ans{$\checkmark$}
    \item \texttt{int a = 012 + 87;} \\
          \texttt{float b = 12.8e+10 ;} \ans{$\checkmark$}
    \item 一个 \texttt{int} 指针变量和一个 \texttt{int} 数组指针变量所占用的存储空间是一样的 \ans{$\checkmark$}
    \item \texttt{int **p, *pa[5];} \\
          \texttt{p = pa ;} \ans{$\checkmark$}
    \item \texttt{int a[5][5];} \\
          \texttt{int (*pa)[5]=a++;} \ans{$\times$}
    \item \texttt{char *str[5];} \\
          \texttt{scanf("\%s",str[0]);} \ans{$\times$}
    \item \texttt{int x, *p =\&x , a[5][5];} \\
          \texttt{*p = a[0][0] + 10 ;} \ans{$\checkmark$}
    \item \texttt{int a = (1e5, 0x1a, 0187);} \ans{$\times$}
    \item \texttt{char x = '\textbackslash 177' ;} \ans{$\checkmark$}
    \item \texttt{char *str = "\textbackslash\textbackslash\textbackslash\textbackslash\textbackslash\textbackslash\textbackslash text\textbackslash";} \ans{$\times$}
\end{enumerate}

\section*{二、计算下列表达式的值（10 分）}

\noindent 假设 \texttt{int} 和 \texttt{unsigned} 类型均为 16 位长度，且各题彼此独立。\\
设：
\begin{lstlisting}
unsigned int a=1, b=2;
int c=4, d=6;
float f;
char str[]= "HUST";
char s1[10]={'A', 'I', 'A'};
char s2[20]={'h', 'u', 's', 't', '\0', 'A', 'I', 'A'};
\end{lstlisting}

\begin{enumerate}[label=(\arabic*), leftmargin=2.2em, itemsep=4pt]
    \item \texttt{f = a>b ? a/b : d/c;} \ans{\texttt{1.0}}
    \item \texttt{!a \&\& --b || c/d;} \ans{\texttt{0}}
    \item \texttt{f = (a - c)/2} \ans{\texttt{32766.0}}
    \item \texttt{f = sizeof(str);} \ans{\texttt{5.0}}
    \item \texttt{f = strlen(s1) + strlen(s2);} \ans{\texttt{7.0}}
\end{enumerate}

\section*{三、程序改错（10 分）}

\noindent 要求：不得改变程序框架，不得重写程序，无需文字说明，直接在代码上添加、删除和修改。

\medskip
\noindent 1、如下程序，删除给定字符串中的数字字符，主函数中输入和输出，调用函数完成删除功能。

\noindent\textbf{修改点：}
\begin{itemize}[leftmargin=2em, itemsep=1pt, topsep=2pt]
    \item 缺少以下两行：\texttt{\#include <string.h>} 和 \texttt{char *delnum(char *s);}
    \item \texttt{char *string[80];} $\longrightarrow$ \texttt{char string[80], *pstr;}
    \item \texttt{char delnum(char **s)} $\longrightarrow$ \texttt{char *delnum(char *s)}
    \item \texttt{char *temp = *s;} $\longrightarrow$ \texttt{char *temp = s;}
    \item \texttt{for (i = 0; s[i] != '\textbackslash 0'; i++)} $\longrightarrow$ \texttt{for (i = 0; s[i] != '\textbackslash 0'; )}
    \item \texttt{strcpy(s+1, s);} $\longrightarrow$ \texttt{strcpy(s + i, s + i + 1);}
\end{itemize}
\begin{lstlisting}
#include <stdio.h>
#include <string.h>
char *delnum(char *s);

void main()
{
    char string[80], *pstr;
    printf("input string:\n");
    gets(string);
    pstr=delnum(string);
    puts(pstr);
}
char *delnum(char *s)
{
    int i;
    char *temp = s;
    for (i = 0; s[i] != '\0'; )
    {
        if (s[i] >= '0' && s[i] <= '9')
            strcpy(s + i, s + i + 1);
        else
            i++;
    }
    return temp;
}
\end{lstlisting}

\medskip
\noindent 2、求 3 个浮点数的平方和与立方和

\noindent\textbf{修改点：}
\begin{itemize}[leftmargin=2em, itemsep=1pt, topsep=2pt]
    \item 缺少函数声明：\texttt{float calc(float x, float y, float z, float *s\_v);}
    \item \texttt{main()} $\longrightarrow$ \texttt{void main()}
    \item \texttt{float a, b;} $\longrightarrow$ \texttt{float a, b, c;}
    \item \texttt{flaot pingf\_v, lif\_v;} $\longrightarrow$ \texttt{float pingf\_v, lif\_v;}
    \item \texttt{scanf("\%f\%f\%f", a, b, c);} $\longrightarrow$ \texttt{scanf("\%f\%f\%f", \&a, \&b, \&c);}
    \item \texttt{pingf\_v = calc(a, b, c, lif\_v);} $\longrightarrow$ \texttt{pingf\_v = calc(a, b, c, \&lif\_v);}
    \item \texttt{calc(float x, float y, float s\_v)} $\longrightarrow$ \\
          \texttt{float calc(float x, float y, float z, float *s\_v)}
    \item \texttt{float *temp;} $\longrightarrow$ \texttt{float temp;}
    \item \texttt{pingf\_v = a*a+b*b+c*c;} $\longrightarrow$ \texttt{temp = x*x+y*y+z*z;}
    \item \texttt{s\_v = a*a*a + b*b*b+c*c*c;} $\longrightarrow$ \texttt{*s\_v = x*x*x + y*y*y+z*z*z;}
    \item \texttt{*temp=pingf\_v} $\longrightarrow$ 去掉
    \item \texttt{return *temp;} $\longrightarrow$ \texttt{return temp;}
\end{itemize}
\begin{lstlisting}
#include <stdio.h>
float calc(float x, float y, float z, float *s_v);
void main()
{
    float a, b, c;
    float pingf_v, lif_v;
    scanf("%f%f%f", &a, &b, &c);
    pingf_v = calc(a, b, c, &lif_v);
    printf("pingf_v=%f, lif_v=%f", pingf_v, lif_v);
}
float calc(float x, float y, float z, float *s_v)
{
    float temp;
    temp = x*x + y*y + z*z;
    *s_v = x*x*x + y*y*y + z*z*z;
    return temp;
}
\end{lstlisting}

\section*{四、程序填空（10 分）}

\noindent 1）编一程序实现一个最简单的计算器的功能，如输入 \texttt{3+5} 回车显示 \texttt{3+5=8}；输错就退出（输入的不是加减乘除的运算就算错）。
\begin{lstlisting}[escapeinside={(*@}{@*)}]
#include <stdio.h>
#include <math.h>
void main( )
{
    (*@\fillans{float a, b}@*) ;
    char op ;
    while((*@\fillans{1}@*) )
    {
        scanf("%f %c %f", &a,&op, &b); //
        if((*@\fillans{(op!='+') \&\& (op!='-') \&\& (op!='*') \&\& (op!='/')}@*) )
        (*@\fillans{break}@*);
        switch(op)
        {
            case '+': printf("%f+%f=%f",a,b,a+b);
                      break;
            case '-': printf("%f-%f=%f",a, b, a-b);
                      break;
            case '*': printf("%f * %f=%f",a,b,a*b);
                      break;
            case '/': if((*@\fillans{fabs(b)<1e-6}@*))
                          printf("除法错");
                      else
                          printf("%f / %f = %f",a,b,a/b);
                      break;
        }
    }
}
\end{lstlisting}

\medskip
\noindent 2）输入一行字符串，统计其中 26 个字母各自出现的次数（不分大小写）。
\begin{lstlisting}[escapeinside={(*@}{@*)}]
#include <string.h>
#include <stdio.h>
void main( )
{
    char str[256];
    unsigned int i;
    static int count[26];
    gets(str);
    for( i=0; str[i]!='\0' ; i++)
    {
        if( (*@\fillans{str[i]>='a' \&\& str[i]<='z'}@*) )
        {
            count[(*@\fillans{str[i]-'a'}@*)]++;
        }
        else if ( (*@\fillans{str[i]>='A' \&\& str[i]<='Z'}@*) )
        {
            (*@\fillans{count[str[i]-'A']++}@*) ;
        }
    }

    for (i=0;i<26; i++)
    printf("字母：%c 的个数=%d\n", (*@\fillans{'a'+i, count[i]}@*));
}
\end{lstlisting}

\section*{五、写出程序运行结果（25 分，结果写在题目的右边）}

\begin{enumerate}[label=\arabic*、, leftmargin=2em, itemsep=8pt]
\item
\begin{lstlisting}
#include<stdio.h>
void main( )
{
    int a=3,b=4;
    int c=5,d=6;
    float f;
    f = (a++,b++,a+b);
    printf("f=%f\n",f);
    f = b/a + d/c ;
    printf("f=%f\n",f);
    f = (a|b)^(c&d);
    printf("f=%f\n",f);
}
\end{lstlisting}
\noindent\textbf{运行结果：}
\begin{verbatim}
f=9.000000
f=2.000000
f=1.000000
\end{verbatim}

\item
\begin{lstlisting}
#include<stdio.h>
void fun( )
{
    int at= 0;
    static int st=0;
    printf("at=%d,  st=%d  st+at=%d\n",at++,st++,st+at);
}
void main( )
{
    int i;
    for(i=0;i<5;i++)
    {
        fun( );
    }
}
\end{lstlisting}
\noindent\textbf{运行结果：}
\begin{verbatim}
at=0,  st=0  st+at=0
at=0,  st=1  st+at=1
at=0,  st=2  st+at=2
at=0,  st=3  st+at=3
at=0,  st=4  st+at=4
\end{verbatim}

\item
\begin{lstlisting}
#include<stdio.h>
void main( )
{
    char str[][20]={ "china",
        "hubei",
        "wuhan",
        "hust-aia" ,
        "c language" };
    int i ;
    for ( i=0 ; i < 5 ; i++)
    {
    printf("%s\n", str[4-i]+i);
    }
}
\end{lstlisting}
\noindent\textbf{运行结果：}
\begin{verbatim}
c language
ust-aia
han
ei
a
\end{verbatim}

\item
\begin{lstlisting}
#include<stdio.h>
void main( )
{
    char *str[]={"china", "hubei", "wuhan", "hust-aia"};
    char **p[] = {str+3 , str+2 , str+1 , str};
    char ***pp = p ;
    printf("%s---", **++pp);
    printf("---%s---", *pp[1]);
    printf("---%s", *--*++pp+2);
}
\end{lstlisting}
\noindent\textbf{运行结果：}
\begin{verbatim}
wuhan------hubei------ina
\end{verbatim}

\newpage
\item
\begin{lstlisting}
#include<stdio.h>
void main()
{
    int a[5][5];
    int i, j;
    for (i = 0; i < 5; i++)
    {
        a[i][0] = 1;
        a[i][i] = 1;
    }
    for (i = 2; i < 5; i++)
    {
        for (j = 1; j < i; j++)
        {
            a[i][j] = a[i - 1][j - 1] + a[i - 1][j];
        }
    }
    for (i = 0; i < 5; i++)
    {
        for (j = 0; j <= i; j++)
        {
            printf("%5d", a[i][j]);
        }
        printf("\n");
    }
}
\end{lstlisting}
\noindent\textbf{运行结果：}
\begin{verbatim}
    1
    1    1
    1    2    1
    1    3    3    1
    1    4    6    4    1
\end{verbatim}
\end{enumerate}

\section*{六、编程（35 分）注意：不得使用全局变量，注意程序结构}

\noindent 1）编程，输入一个数值 $x$，计算 $y$ 的结果。（8~分）
\[
y = \begin{cases}
x+1 & 0 \le x < 2 \\
2x+2 & 2 \le x < 4 \\
3x+3 & 4 \le x < 6 \\
4x+4 & 6 \le x < 8
\end{cases}
\]
\begin{lstlisting}
#include <stdio.h>
void main( )
{
    float x, y;
    printf("input the number x=");
    scanf("%f",&x);
    if( x >= 0 && x < 2 )
    {
        y = x + 1;
    }
    else if (x >= 2 && x < 4)
    {
        y = 2*x + 2;
    }
    else if ( x >= 4 && x < 6)
    {
        y = 3*x + 3 ;
    }
    else if ( x >= 6 && x < 8)
    {
        y = 4*x + 4;
    }
    else
    {
        printf("error in input data\n");
        return;
    }
    printf("y=%.2f",y);
}
\end{lstlisting}

\medskip
\noindent 2、编程求序列：$\dfrac{1}{2}, \dfrac{3}{4}, \dfrac{5}{8}, \dfrac{7}{16}, \dfrac{9}{32} \cdots\cdots$，所有大于等于 $0.000001$ 的数据项之和，显示输出计算的结果。（8~分）
\begin{lstlisting}
#include <stdio.h>
#include <math.h>
void main( )
{
    float sum, a, b; //sum 代表和, a 为分子, b 为分母
    sum = 0;
    a = 1;           // 分子赋初值
    b = 2;           // 分母赋初值
    while (a/b >= 1e-6)
    {
        sum = sum + a/b; // 累加一项
        a = a + 2;       // 求下一项的分子
        b = b * 2;       // 求下一项的分母
    }
    printf(" sum = %f", sum);
}
\end{lstlisting}

\medskip
\noindent 3、编程序，要求主函数中输入 $N$ 行英文，被调用的函数返回其中最长的那个单词（如有多个，返回最后出现的）并在主函数中输出。（10~分）
\begin{lstlisting}
#define N 10
#include<string.h>
#include<stdio.h>
char *findstr(char (*pstr)[80]);
void main()
{
    int i;
    char str[N][80], *ps;
    for(i=0; i<N; i++)
    {
        gets(str[i]);
    }
    ps = findstr(str);
    puts(ps);
}
char *findstr(char (*pstr)[80])
{
    int i, j, k=0;
    static char maxstr[80], tempstr[80];
    for(i=0; i<N; i++)
    {
        for(j=0; pstr[i][j]!='\0'; j++)
        {
            if((pstr[i][j]>='a' && pstr[i][j]<='z')
               || (pstr[i][j]>='A' && pstr[i][j]<='Z'))
            {
                tempstr[k++]=pstr[i][j];
            }
            else
            {
                tempstr[k]='\0';
                if(strlen(tempstr)>=strlen(maxstr))
                    strcpy(maxstr,tempstr);
                k=0;
            }
        }
        tempstr[k]='\0';
        if(strlen(tempstr)>=strlen(maxstr))
            strcpy(maxstr,tempstr);
        k=0;
    }
    return maxstr;
}
\end{lstlisting}

\medskip
\noindent 4、某班有学生 20 名参加期末考试，每名学生信息由姓名、学号、性别、外语、数学、C 课设和总分组成，试编程要求：（9~分）
\begin{enumerate}[label=(\arabic*), leftmargin=2.2em, itemsep=2pt]
    \item 学生信息由键盘输入；
    \item 能按 C 课设成绩查找学生信息；
    \item 按总分高低顺序输出 C 课设挂科的男生信息。
\end{enumerate}
\noindent （每小题写一函数，通过 \texttt{main} 函数调用实现）
\begin{lstlisting}
struct student{
    char name[10];
    char no[20];
    char sex;
    float engl;
    float math;
    float Cpro;
    float sum;
};
#define N 20
#include <stdio.h>
void input(struct student *p);
void findstu(struct student *p);
void outputstu(struct student *p);

void main( )
{
    struct student stu[N];
    input(stu);
    findstu(stu);
    outputstu(stu);
}
void input(struct student *p)
{
    int i;
    for(i=0; i<N; i++, p++)
    {
        scanf("%s%s %c%f%f%f", p->name, p->no, &p->sex,
              &p->engl, &p->math, &p->Cpro);
        p->sum = p->engl + p->math + p->Cpro;
    }
}
void findstu(struct student *p)
{
    int i; float x;
    while(1)
    {
        printf("请输入 C 课设成绩"); // 系统设计成输入 0-100 之外的数退出,
        // 输入正确的 C 课设分数 x, 则显示分数大于等于 x 的所有学生
        scanf("%f", &x);
        if(x<0 || x>100) break;
        for(i=0; i<N; i++)
        {
            if(p[i].Cpro >= x)
                printf("%10s%20s %c%10f%10f%10f%10f\n",
                       p[i].name, p[i].no, p[i].sex,
                       p[i].engl, p[i].math, p[i].Cpro, p[i].sum);
        }
    }
}
void outputstu(struct student *p)
{
    int i, j;
    struct student temp;
    for(i=0; i<N-1; i++)
        for(j=0; j<N-1-i; j++)
        {
            if(p[j].sum < p[j+1].sum)
            {
                temp = p[j];
                p[j] = p[j+1];
                p[j+1] = temp;
            }
        }
    for(i=0; i<N; i++)
    {
        if(p[i].Cpro<60 && p[i].sex=='M')
            printf("%10s%20s %c%10f%10f%10f%10f\n",
                   p[i].name, p[i].no, p[i].sex,
                   p[i].engl, p[i].math, p[i].Cpro, p[i].sum);
    }
}
\end{lstlisting}

\end{document}
